\documentclass[10pt,a4paper]{amsart}
\usepackage[margin=19mm]{geometry}
\usepackage{amsmath,amssymb,mathtools}
\usepackage[scr=rsfs,cal=euler]{mathalfa}
\usepackage{xfrac}
\usepackage[unicode,hidelinks,bookmarks=false]{hyperref}
\newcommand{\ie}{i.e.\ }
\newcommand{\cf}{cf.\ }
\newcommand{\resp}{respectively\ }
\newcommand{\Subsection}{\subsection}
\providecommand{\nobreakdash}{\nobreak\hbox{-}}
\setlength{\emergencystretch}{2em}
\multlinegap0pt
\usepackage{bm}
\usepackage{bbm}
\usepackage{dsfont}
\usepackage{xspace}
\usepackage{cancel}
\usepackage{mathtools}
\usepackage{amsthm}
\newtheorem{theorem}{Theorem}
\newtheorem{lemma}[theorem]{Lemma}
\newtheorem{remark}[theorem]{Remark}
\newcommand{\Fr}{\operatorname{Fr}}
\newcommand{\LIP}{\operatorname{LIP}}
\newcommand{\Mod}{\operatorname{Mod}}
\newcommand{\N}{\mathbb{N}}
\newcommand{\dom}{\operatorname{dom}}
\newcommand{\gap}{\operatorname{gap}}
\newcommand{\ud}{\mathrm {d}}
\title[Fragment-wise differentiable structures]{Fragment-wise differentiable structures}
\author{David Bate, Sylvester Eriksson-Bique, Elefterios Soultanis}
\date{Source: arXiv:2402.11284v2 (CC BY 4.0)}
\begin{document}
\maketitle

\noindent\emph{Source and licence.} Fragment-wise differentiable structures. Version arXiv:2402.11284v2 (CC BY 4.0), distributed under the Creative Commons Attribution 4.0 International licence, as recorded in the arXiv licence metadata for this exact version and shown on the abstract page https://arxiv.org/abs/2402.11284v2. Full rights notice: licenses/arxiv-2402.11284-CC-BY-4.0.txt.\par\smallskip

\noindent\textit{Editorial scope.} Counted unit: the Lemma whose source label is \texttt{lem:lim-star-ug-weak} in the article (source0.tex lines 605--607, source label \texttt{lem:lim-star-ug-weak}), delivered together with the article's own complete original proof of it (source0.tex lines 608--637, one \texttt{proof} environment of 30 lines). No mathematics is written, repaired or re-derived: statement and proof are the article's own text line for line. Required context retained uncounted: rmk:uppergradchar (lines 547--553); lem:lim-star-ug (lines 585--587). Every definition, display and label of those lines is retained unchanged. Omitted from this package and not redistributed: 1--546, 554--584, 588--604, 638--1718. Nothing inside a retained excerpt is omitted or altered: every retained body is delivered exactly as the article's lines are written, so no omission receipt of any kind accompanies this package. Citations: the retained excerpts carry no citation command at all, so no bibliography entry of the article is transcribed and no reference is redistributed. Reference adaptation: none was needed and none was made; every reference inside the delivered unit resolves inside it, so no unresolved reference or citation is produced. Printed slips kept literal: no printed word, symbol, hypothesis or number of the counted statement or of its proof was altered. Layout and class: the article's own class is \texttt{amsart} with packages including amsmath, amsthm, amsfonts, amssymb, graphicx, tikz-cd, bm, comment, hyperref, todonotes, mathrsfs; the wrapper is the lane's pinned \texttt{amsart} editorial wrapper at 10pt on A4, which restates the theorem-like environments the retained text needs and adds the article's own macro definitions that the retained text needs (\texttt{\textbackslash Fr}, \texttt{\textbackslash LIP}, \texttt{\textbackslash Mod}, \texttt{\textbackslash N}, \texttt{\textbackslash dom}, \texttt{\textbackslash gap}, \texttt{\textbackslash ud}), with extra wrapper packages (bm, bbm, dsfont, xspace, cancel, mathtools). The delivered numbering is the unit's own. Assets: no retained span contains a figure environment, a graphics-inclusion command, a PDF-page inclusion, a table or a TikZ picture; no image file is copied into this package, the package declares no asset dependency and no image file is redistributed with it. Licence: version arXiv:2402.11284v2 of this article is distributed under the Creative Commons Attribution 4.0 International licence, as recorded in the arXiv licence metadata for this exact version and shown on the abstract page https://arxiv.org/abs/2402.11284v2. A full rights notice is delivered at licenses/arxiv-2402.11284-CC-BY-4.0.txt. Characters: every retained excerpt is pure ASCII, so the delivered engine, XeLaTeX with its default Latin Modern text font, reports no missing character.
\begin{remark}\label{rmk:uppergradchar}
In the proof below we use the observation that a Borel function $\rho\in L^\infty(\mu)$ is a weak $\ast$-upper gradient of $f\in \LIP(X)$ if and only if 
\begin{align*}
|(f\circ\gamma)'(t)|\le \rho(\gamma_t)|\gamma'_t|\quad a.e.\ t\in \operatorname{dom}(\gamma)
\end{align*}
for $\Mod_\infty$-a.e. $\gamma\in \Fr(X)$. The proof also uses the concept of a measure theoretic infimum, which is also sometimes referred to as a lattice infimum.
\end{remark} 

\begin{lemma}\label{lem:lim-star-ug}
Suppose $(f_j)\subset \LIP(X)$ converges pointwise to $f\in \LIP(X)$ and $\sup_j\LIP(f_j)<\infty$. If $\rho_j\in L^\infty(\mu)$ is a weak $\ast$-upper gradient of $f_j$ for each $j$ and $(\rho_j)$ converges pointwise $\mu$-a.e. to $\rho\in L^\infty$ with $\sup_j\|\rho_j\|_{L^\infty(\mu)}<\infty$, then $\rho$ is a weak $\ast$-upper gradient of $f$.
\end{lemma}

\begin{lemma}\label{lem:lim-star-ug-weak}
Suppose $(f_j)\subset \LIP(X)$ converges pointwise to $f\in \LIP(X)$ and $\sup_j\LIP(f_j)<\infty$. If $\rho_j\in L^\infty(\mu)$ is a weak $\ast$-upper gradient of $f_j$ for each $j$ and $(\rho_j)$ $w^\ast$-converges in $L^\infty(\mu)$ to $\rho$, then $\rho$ is a weak $\ast$-upper gradient of $f$.
\end{lemma}

\begin{proof}
First, by Banach-Steinhaus, we get $\sup_j\|\rho_j\|_{L^\infty(\mu)}<\infty$.
We will show that there exists non-negative numbers $\alpha_k(n), \dots \alpha_{N_n}(n)$, $N_n\in \N$ (with $n\leq N_n)$ and with $\sum_{j=n}^{N_n} \alpha_j(n) = 1$, so that
\[
\tilde{\rho}_n=\sum_{j=n}^{N_n} \alpha_j(n) \rho_j \stackrel{n\to\infty}{\longrightarrow} \rho
\]
converges pointwise almost everywhere and with $\sup_j\|\tilde{\rho}_j\|_{L^\infty(\mu)}<\infty$. If this is true, then it is direct to verify that $\tilde{\rho}_k$ is a weak $\ast$-upper gradient for $\tilde{f}_n=\sum_{j=n}^{N_n} \alpha_j(n) f_j$. Indeed, for $\Mod_\infty$-a.e. curve fragment $\gamma$, and any $s,t\in \dom(\gamma)$, we have
\begin{align*}
|\tilde{f}_n(\gamma_t)-\tilde{f}_n(\gamma_s)|&\le \sum_{j=n}^{N_n} \alpha_j(n) |f_j(\gamma_t)-f_j(\gamma_s)|
\\
& \leq \sum_{j=n}^{N_n} \alpha_j(n) \bigg(\int_{\gamma}\rho_j\ud s+\LIP(f_j)\gap(\gamma)\bigg)\\
\leq & \int_\gamma\tilde{\rho}_n \ud s+C\gap(\gamma),
\end{align*}
where $C=\sup_j \LIP(f_j)$. From this and Remark \ref{rmk:uppergradchar}, similar to the previous lemma, we get that $\tilde{\rho}_n$ is a weak $\ast$-upper gradient for $\tilde{f}_n$. Now, the claim follows from Lemma \ref{lem:lim-star-ug}. We are left to prove the existence of $\alpha_j(n)$ and $N_n$ such that the previous claims hold.

Fix $n\in \N$ and fix $x_0\in X$. Since $L^\infty(B(x_0,n),\mu) \subset L^2(B(x_0,n),\mu)\newline \subset L^1(B(x_0,n),\mu)$, we have that the sequence $\rho_j \chi_R \to \rho \chi_R$ weakly in $L^2(B(x_0,n),\mu)$. Thus, by Mazur's lemma, together with passing to a subsequence, we can find non-negative values $\beta^n_k(k), \dots \beta^n_{M^n_k}(k)$, $k\leq M^n_k\in \N$, with $\sum_{j=k}^{M^n_k} \beta_j(k) = 1$, so that
\[
\tilde{\rho}^n_k=\sum_{j=k}^{M^n_k} \beta_j(k) \rho_j \stackrel{k\to\infty}{\longrightarrow} \rho
\]
converges almost everywhere in $B(x_0,n)$. So, for some index $m_n\in \N$ we have $m_n\geq n$ and $\mu(\{x\in B(x_0,n): |\tilde{\rho}^n_{m_n}(x)-\rho(x)|\geq\frac{1}{n}\}) \leq 2^{-n}$. Set $N_n=M_{m_n}$, and $\alpha_j(n)=\beta_j(m)$ for $j\in [m,M_m]$, and otherwise $\alpha_j(n)=0$ for $j\in [n,m]$.

By the following Borel-Cantelli-type argument we have that almost everywhere
\[
\tilde{\rho}_n=\sum_{j=n}^{N_n} \alpha_j(n) \rho_j \stackrel{n\to\infty}{\longrightarrow} \rho.
\]
Indeed, let $A_n=\{x\in B(x_0,n): |\tilde{\rho}^n_{m_n}(x)-\rho(x)|\geq\frac{1}{n}\}=\{\{x\in B(x_0,n): |\tilde{\rho}_{n}(x)-\rho(x)|\geq\frac{1}{n}\}.$ Then for the limit superior set $\displaystyle \limsup_{n\to \infty} A_n = \bigcap_{m=1}^\infty \bigcup_{n\geq m} A_n$  we have 
\[
\mu\left(\bigcap_{m=1}^\infty \bigcup_{n\geq m} \limsup_{n\to \infty} A_n\right) = 0,
\]and for every $x\not\in \limsup_{n\to \infty} A_n$, we get pointwise convergence $\lim_{n\to \infty}\tilde{\rho}_n(x)=\rho(x)$. 
\end{proof}

\end{document}
